\documentclass[UTF8]{ctexart}
\usepackage{amsmath}
\usepackage{amssymb}
\usepackage{geometry}
\usepackage{xcolor}
\usepackage{amsthm}
\usepackage{tcolorbox}
\usepackage{hyperref}
\usepackage{titlesec}

\geometry{a4paper, scale=0.8}

% --- 颜色定义 ---
\definecolor{myblue}{RGB}{30,144,255}
\definecolor{mygreen}{RGB}{60,179,113}
\definecolor{myorange}{RGB}{255,165,0}
\definecolor{boxback}{RGB}{245,245,245}

% --- 标题样式 ---
\titleformat{\section}{\Large\bfseries\color{myblue}}{\thesection}{1em}{}
\titleformat{\subsection}{\large\bfseries\color{mygreen}}{\thesubsection}{1em}{}
\titleformat{\subsubsection}{\bfseries\color{myorange}}{\thesubsubsection}{1em}{}

% --- 定理环境 ---
\newtheoremstyle{mydef}
  {} % Space above
  {} % Space below
  {} % Body font
  {} % Indent amount
  {\bfseries} % Thm head font
  {.} % Punctuation after thm head
  {.5em} % Space after thm head
  {} % Thm head spec

\theoremstyle{mydef}
\newtheorem{definition}{定义}[section]
\newtheorem{property}{性质}[section]
\newtheorem{theorem}{定理}[section]
\newtheorem{conclusion}{结论}[section]

% --- tcolorbox 环境 ---
\newtcolorbox{mybox}[2][]{
  colback=boxback,
  colframe=mygreen!75!black,
  fonttitle=\bfseries,
  coltitle=black,
  title=#2,
  #1
}

% --- 超链接设置 ---
\hypersetup{
    colorlinks=true,
    linkcolor=myblue,
    filecolor=magenta,      
    urlcolor=cyan,
    citecolor=mygreen,
    pdftitle={线性代数知识点总结},
    pdfauthor={GitHub Copilot},
}

\title{线性代数知识点总结}
\author{GitHub Copilot}
\date{\today}

\begin{document}

\maketitle
\tableofcontents
\newpage

\section{行列式的定义}

\subsection{n 级排列}

\begin{enumerate}
    \item \textbf{定义}：由 $1, 2, \cdots, n$ 组成的一个有序数组称为一个 $n$ 级排列，通常记为 $i_1 i_2 \cdots i_n$。
    \item \textbf{逆序}：在一个排列中，如果一对数的前后位置与大小顺序相反，即前面的数大于后面的数，则它们称为一个逆序。
    \item \textbf{逆序数}：一个排列中的逆序总数，通常记为 $\tau(i_1 i_2 \cdots i_n)$。
    \item \textbf{奇（偶）排列}：逆序数为奇数（偶数）的排列。
    \item \textbf{对换}：把一个排列中某两个数的位置互换，而其余的数不动，就得到另一个排列，这样一个变换称为一个对换。
\end{enumerate}

\subsection{n 级排列的性质}

\begin{enumerate}
    \item 任意一个排列经过一个对换后，奇偶性改变。
    \item $n$ 级排列共有 $n!$ 种，奇偶排列各占一半。
\end{enumerate}

\subsection{n 阶行列式}

$n$ 阶行列式
$$
\begin{vmatrix}
a_{11} & a_{12} & \cdots & a_{1n} \\
a_{21} & a_{22} & \cdots & a_{2n} \\
\vdots & \vdots & \vdots & \vdots \\
a_{n1} & a_{n2} & \cdots & a_{nn}
\end{vmatrix}
$$
是所有取自不同行不同列的 $n$ 个元素的乘积 $a_{1j_1} a_{2j_2} \cdots a_{nj_n}$ 的代数和，这里 $j_1 j_2 \cdots j_n$ 是一个 $n$ 级排列。

当 $j_1 j_2 \cdots j_n$ 是偶排列时，该项前面带正号；
当 $j_1 j_2 \cdots j_n$ 是奇排列时，该项前面带负号。

即：
$$
\begin{vmatrix}
a_{11} & a_{12} & \cdots & a_{1n} \\
a_{21} & a_{22} & \cdots & a_{2n} \\
\vdots & \vdots & \vdots & \vdots \\
a_{n1} & a_{n2} & \cdots & a_{nn}
\end{vmatrix}
= \sum_{j_1 j_2 \cdots j_n} (-1)^{\tau(j_1 j_2 \cdots j_n)} a_{1j_1} a_{2j_2} \cdots a_{nj_n}
$$
其中 $\sum_{j_1 j_2 \cdots j_n}$ 表示对所有 $n$ 级排列求和。

$n$ 阶行列式有时简记为 $|a_{ij}|_n$，而且有如下另外两种类似的定义：
$$
|a_{ij}|_n = \sum_{i_1 i_2 \cdots i_n} (-1)^{\tau(i_1 i_2 \cdots i_n)} a_{i_1 1} a_{i_2 2} \cdots a_{i_n n}
$$
和
$$
|a_{ij}|_n = \sum_{\substack{i_1 i_2 \cdots i_n \\ j_1 j_2 \cdots j_n}} (-1)^{\tau(i_1 i_2 \cdots i_n) + \tau(j_1 j_2 \cdots j_n)} a_{i_1 j_1} a_{i_2 j_2} \cdots a_{i_n j_n}
$$
由 $n$ 级排列的性质可知，$n$ 阶行列式共有 $n!$ 项，其中冠以正号的项和冠以负号的项（不算元素本身所带的负号）各占一半。

\section{常见行列式}

\subsection{二阶行列式}
$$
\begin{vmatrix}
a_{11} & a_{12} \\
a_{21} & a_{22}
\end{vmatrix}
= a_{11}a_{22} - a_{12}a_{21}
$$

\subsection{三阶行列式}
$$
\begin{vmatrix}
a_{11} & a_{12} & a_{13} \\
a_{21} & a_{22} & a_{23} \\
a_{31} & a_{32} & a_{33}
\end{vmatrix}
= a_{11}a_{22}a_{33} + a_{12}a_{23}a_{31} + a_{13}a_{21}a_{32} - a_{13}a_{22}a_{31} - a_{12}a_{21}a_{33} - a_{11}a_{23}a_{32}
$$

\subsection{上三角、下三角、对角行列式}
$$
\begin{vmatrix}
a_{11} & * & \cdots & * \\
0 & a_{22} & \cdots & * \\
\vdots & \vdots & \ddots & \vdots \\
0 & 0 & \cdots & a_{nn}
\end{vmatrix}
=
\begin{vmatrix}
a_{11} & 0 & \cdots & 0 \\
* & a_{22} & \cdots & 0 \\
\vdots & \vdots & \ddots & \vdots \\
* & * & \cdots & a_{nn}
\end{vmatrix}
=
\begin{vmatrix}
a_{11} & 0 & \cdots & 0 \\
0 & a_{22} & \cdots & 0 \\
\vdots & \vdots & \ddots & \vdots \\
0 & 0 & \cdots & a_{nn}
\end{vmatrix}
= a_{11} a_{22} \cdots a_{nn}
$$

\subsection{副对角线方向的行列式}
$$
\begin{vmatrix}
* & \cdots & * & a_{1n} \\
* & \cdots & a_{2, n-1} & 0 \\
\vdots & \dots & \vdots & \vdots \\
a_{n1} & 0 & \cdots & 0
\end{vmatrix}
= (-1)^{\frac{n(n-1)}{2}} a_{1n} a_{2, n-1} \cdots a_{n1}
$$
(注：图中给出的公式形式针对的是从右上到左下的副对角线形式，且左上部分为0，或者右下部分为0的情况，结果符号由排列逆序数决定，即 $\frac{n(n-1)}{2}$)

\subsection{范德蒙 (Vandermonde) 行列式}
$$
\begin{vmatrix}
1 & 1 & 1 & \cdots & 1 \\
a_1 & a_2 & a_3 & \cdots & a_n \\
a_1^2 & a_2^2 & a_3^2 & \cdots & a_n^2 \\
\vdots & \vdots & \vdots & & \vdots \\
a_1^{n-1} & a_2^{n-1} & a_3^{n-1} & \cdots & a_n^{n-1}
\end{vmatrix}
= \prod_{1 \le j < i \le n} (a_i - a_j)
$$


\section{行列式的性质}

\begin{enumerate}
    \item \textbf{性质 1}：行列式 $D$ 与它的转置行列式 $D^T$ 相等。
    \item \textbf{性质 2}：互换行列式的两行（列），行列式变号（交换 $r_i$ 行和 $r_j$ 行，记为 $r_i \leftrightarrow r_j$；交换 $c_i$ 列和 $c_j$ 列，记为 $c_i \leftrightarrow c_j$）。
    \item \textbf{性质 3}：行列式的某一行（列）中所有的元素都乘以同一数 $k$，等于用数 $k$ 乘此行列式。
    \begin{itemize}
        \item \textbf{推论}：行列式中某一行（列）的所有元素的公因子可以提到整个行列式的外面。
    \end{itemize}
    \item \textbf{性质 4}：行列式中如果有两行（列）元素成比例，则此行列式等于零。
    \item \textbf{性质 5}：若行列式的某一列（行）的所有元素都是两数之和，例如第 $i$ 列的元素都是两数之和，即
    $$
    D = \begin{vmatrix}
    a_{11} & a_{12} & \cdots & (a_{1i} + a'_{1i}) & \cdots & a_{1n} \\
    a_{21} & a_{22} & \cdots & (a_{2i} + a'_{2i}) & \cdots & a_{2n} \\
    \vdots & \vdots & & \vdots & & \vdots \\
    a_{n1} & a_{n2} & \cdots & (a_{ni} + a'_{ni}) & \cdots & a_{nn}
    \end{vmatrix}
    $$
    则 $D$ 等于下列两个行列式之和：
    $$
    D = \begin{vmatrix}
    a_{11} & a_{12} & \cdots & a_{1i} & \cdots & a_{1n} \\
    a_{21} & a_{22} & \cdots & a_{2i} & \cdots & a_{2n} \\
    \vdots & \vdots & & \vdots & & \vdots \\
    a_{n1} & a_{n2} & \cdots & a_{ni} & \cdots & a_{nn}
    \end{vmatrix}
    +
    \begin{vmatrix}
    a_{11} & a_{12} & \cdots & a'_{1i} & \cdots & a_{1n} \\
    a_{21} & a_{22} & \cdots & a'_{2i} & \cdots & a_{2n} \\
    \vdots & \vdots & & \vdots & & \vdots \\
    a_{n1} & a_{n2} & \cdots & a'_{ni} & \cdots & a_{nn}
    \end{vmatrix}
    $$
    \item \textbf{性质 6}：把行列式的某一行（列）的各元素乘以同一数然后加到另一行（列）对应的元素上，行列式不变。
    例如以数 $k$ 乘第 $j$ 列加到第 $i$ 列上（记作 $c_i + k c_j$），有
    $$
    \begin{vmatrix}
    a_{11} & \cdots & a_{1i} & \cdots & a_{1j} & \cdots & a_{1n} \\
    a_{21} & \cdots & a_{2i} & \cdots & a_{2j} & \cdots & a_{2n} \\
    \vdots & & \vdots & & \vdots & & \vdots \\
    a_{n1} & \cdots & a_{ni} & \cdots & a_{nj} & \cdots & a_{nn}
    \end{vmatrix}
    \xrightarrow{c_i + k c_j}
    \begin{vmatrix}
    a_{11} & \cdots & (a_{1i} + k a_{1j}) & \cdots & a_{1j} & \cdots & a_{1n} \\
    a_{21} & \cdots & (a_{2i} + k a_{2j}) & \cdots & a_{2j} & \cdots & a_{2n} \\
    \vdots & & \vdots & & \vdots & & \vdots \\
    a_{n1} & \cdots & (a_{ni} + k a_{nj}) & \cdots & a_{nj} & \cdots & a_{nn}
    \end{vmatrix}
    \quad (i \neq j)
    $$
    （以数 $k$ 乘第 $j$ 行加到第 $i$ 行上，记作 $r_i + k r_j$）。
\end{enumerate}

\section{行列式按行（列）展开}

\subsection{基本概念}

\begin{enumerate}
    \item \textbf{余子式}：在 $n$ 阶行列式 $D = |a_{ij}|$ 中，去掉元素 $a_{ij}$ 所在的第 $i$ 行和第 $j$ 列后，余下的 $n-1$ 阶行列式，称为 $a_{ij}$ 的余子式，记为 $M_{ij}$。
    \item \textbf{代数余子式}：$A_{ij} = (-1)^{i+j}M_{ij}$ 称为 $a_{ij}$ 的代数余子式。
    \item \textbf{$k$ 阶子式}：在 $n$ 阶行列式 $D = |a_{ij}|$ 中，任意选定 $k$ 行 $k$ 列 ($1 \le k \le n$)，位于这些行列交叉处的 $k^2$ 个元素，按原来顺序构成一个 $k$ 阶行列式，称为 $D$ 的一个 $k$ 阶子式。
\end{enumerate}

\subsection{展开定理}

\begin{enumerate}
    \item \textbf{按一行（列）展开}
    \begin{enumerate}
        \item 行列式等于它的任一行（列）的各元素与其对应的代数余子式乘积之和，即：
        \begin{itemize}
            \item 按第 $i$ 行展开：$D = a_{i1}A_{i1} + a_{i2}A_{i2} + \cdots + a_{in}A_{in} \quad (i = 1, 2, \cdots, n)$
            \item 按第 $j$ 列展开：$D = a_{1j}A_{1j} + a_{2j}A_{2j} + \cdots + a_{nj}A_{nj} \quad (j = 1, 2, \cdots, n)$
        \end{itemize}
        \item 行列式某一行（列）的元素与另一行（列）的对应元素的代数余子式乘积之和等于零，即：
        $$ a_{i1}A_{j1} + a_{i2}A_{j2} + \cdots + a_{in}A_{jn} = 0, \quad i \neq j $$
        或
        $$ a_{1i}A_{1j} + a_{2i}A_{2j} + \cdots + a_{ni}A_{nj} = 0, \quad i \neq j $$
    \end{enumerate}
    \item \textbf{按 $k$ 行（$k$ 列）展开（拉普拉斯定理）}：在 $n$ 阶行列式中，任意取定 $k$ 行（$k$ 列）($1 \le k \le n-1$)，由这 $k$ 行（$k$ 列）组成的所有 $k$ 阶子式与它们的代数余子式的乘积之和等于行列式的值。
\end{enumerate}

\section{行列式的计算}

行列式的计算方法有很多种，大致有以下几种思路：
\begin{enumerate}
    \item 思路 1：利用行列式的定义；
    \item 思路 2：利用行列式的性质；
    \item 思路 3：利用行列式的行列展开。
\end{enumerate}

\section{克莱姆法则}

\subsection{克莱姆法则}
如果含 $n$ 个未知量 $n$ 个方程的线性方程组
$$
\begin{cases}
a_{11}x_1 + a_{12}x_2 + \cdots + a_{1n}x_n = b_1 \\
a_{21}x_1 + a_{22}x_2 + \cdots + a_{2n}x_n = b_2 \\
\cdots \cdots \cdots \cdots \cdots \cdots \\
a_{n1}x_1 + a_{n2}x_2 + \cdots + a_{nn}x_n = b_n
\end{cases}
$$
的系数行列式不等于零，即
$$
D = \begin{vmatrix}
a_{11} & \cdots & a_{1n} \\
\vdots & & \vdots \\
a_{n1} & \cdots & a_{nn}
\end{vmatrix} \neq 0
$$
则方程组有唯一解
$$
x_1 = \frac{D_1}{D}, \quad x_2 = \frac{D_2}{D}, \quad \cdots, \quad x_n = \frac{D_n}{D}
$$
其中 $D_j (j=1, 2, \cdots, n)$ 是把系数行列式 $D$ 中第 $j$ 列的元素用方程组右端的常数项代替后所得到的 $n$ 阶行列式，即
$$
D_j = \begin{vmatrix}
a_{11} & \cdots & a_{1, j-1} & b_1 & a_{1, j+1} & \cdots & a_{1n} \\
\vdots & & \vdots & \vdots & \vdots & & \vdots \\
a_{n1} & \cdots & a_{n, j-1} & b_n & a_{n, j+1} & \cdots & a_{nn}
\end{vmatrix}
$$

\subsection{齐次线性方程组}
含 $n$ 个未知量 $n$ 个方程的齐次线性方程组
$$
\begin{cases}
a_{11}x_1 + a_{12}x_2 + \cdots + a_{1n}x_n = 0 \\
a_{21}x_1 + a_{22}x_2 + \cdots + a_{2n}x_n = 0 \\
\cdots \cdots \cdots \cdots \cdots \cdots \\
a_{n1}x_1 + a_{n2}x_2 + \cdots + a_{nn}x_n = 0
\end{cases}
$$
\begin{itemize}
    \item 只有零解的充分必要条件是系数行列式 $D \neq 0$；
    \item 有非零解的充分必要条件是 $D = 0$。
\end{itemize}


\newpage
\section{矩阵}

\subsection{矩阵的概念}
由 $m \times n$ 个数 $a_{ij} (i=1, 2, \cdots, m; j=1, 2, \cdots, n)$ 按一定次序排成的 $m$ 行 $n$ 列的矩形数表
$$
\begin{bmatrix}
a_{11} & a_{12} & \cdots & a_{1n} \\
a_{21} & a_{22} & \cdots & a_{2n} \\
\vdots & \vdots & & \vdots \\
a_{m1} & a_{m2} & \cdots & a_{mn}
\end{bmatrix}
$$
称为 $m \times n$ 矩阵（$m$ 行 $n$ 列矩阵）。$a_{ij}$ 叫做矩阵的元素。矩阵简记为 $A = (a_{ij})_{m \times n}$ 或 $A = (a_{ij})$。

\begin{itemize}
    \item 当 $m=n$ 时，即矩阵的行数与列数相同时，称 $A$ 为 $n$ 阶方阵；
    \item 当 $m=1$ 时，矩阵只有一行，称为行矩阵，记为 $A = (a_{11}, a_{12}, \cdots, a_{1n})$。这样的行矩阵也称为 $n$ 维行向量；
    \item 当 $n=1$ 时，矩阵只有一列，称为列矩阵，记为 $A = \begin{bmatrix} a_{11} \\ a_{21} \\ \vdots \\ a_{m1} \end{bmatrix}$。这样的列矩阵也称为 $m$ 维列向量。
    \item 矩阵 $A$ 中各元素变号得到的矩阵叫做 $A$ 的负矩阵，记作 $-A$，即 $-A = (-a_{ij})_{m \times n}$。
    \item 如果矩阵 $A$ 的所有元素都是 0，则 $A$ 称为零矩阵，记为 $0$。
\end{itemize}

\subsection{矩阵的运算}
\begin{enumerate}
    \item \textbf{矩阵的相等}：设 $A=(a_{ij})_{m \times n}, B=(b_{ij})_{m \times n}$。如果 $a_{ij} = b_{ij} (i=1, \cdots, m; j=1, \cdots, n)$，则称矩阵 $A$ 与 $B$ 相等，记作 $A=B$。
    \item \textbf{矩阵的加、减法}：设 $A=(a_{ij})_{m \times n}, B=(b_{ij})_{m \times n}, C=(c_{ij})_{m \times n}$。其中 $c_{ij} = a_{ij} \pm b_{ij}$。则称 $C$ 为矩阵 $A$ 与 $B$ 的和（或差），记为 $C = A \pm B$。
    \item \textbf{数与矩阵的乘法}：设 $k$ 为一个常数，$A=(a_{ij})_{m \times n}, C=(c_{ij})_{m \times n}$。其中 $c_{ij} = k a_{ij}$。则称矩阵 $C$ 为数 $k$ 与矩阵 $A$ 的数量乘积，简称数乘，记为 $kA$。
    \item \textbf{矩阵的乘法}：设 $A=(a_{ij})_{m \times s}, B=(b_{ij})_{s \times n}, C=(c_{ij})_{m \times n}$。其中 $c_{ij} = \sum_{k=1}^s a_{ik}b_{kj} \quad (i=1, \cdots, m; j=1, \cdots, n)$。则称矩阵 $C$ 为矩阵 $A$ 与 $B$ 的乘积，记为 $AB$，即 $C=AB$。
    \item \textbf{方阵的幂运算}：对 $n$ 阶方阵 $A$ 定义 $A^k = \underbrace{A \cdot A \cdots A}_{k \text{个}}$，称为 $A$ 的 $k$ 次幂。
    \item \textbf{矩阵的转置}：把矩阵 $A=(a_{ij})_{m \times n}$ 的行行列互换而得到的矩阵 $(a_{ji})_{n \times m}$ 称为 $A$ 的转置矩阵，记为 $A^T$ (或 $A'$）。
    \item \textbf{方阵的行列式}：方阵 $A$ 的元素按原来的位置构成的行列式，称为方阵 $A$ 的行列式，记为 $|A|$。若 $|A|=0$，称 $A$ 为奇异矩阵，否则称为非奇异矩阵。
\end{enumerate}

\subsection{矩阵的运算公式}
\begin{enumerate}
    \item \textbf{关于矩阵的加法运算公式}
    \begin{itemize}
        \item $A + B = B + A$
        \item $(A + B) + C = A + (B + C)$
        \item $A + (-A) = 0$
        \item $(A - B) = A + (-B)$
    \end{itemize}
    \item \textbf{关于数乘运算的公式}
    \begin{itemize}
        \item $1 A = A$
        \item $k(lA) = (kl)A$
        \item $k(A + B) = kA + kB$
        \item $(k + l)A = kA + lA$
    \end{itemize}
    \item \textbf{关于矩阵的乘法运算的公式}
    \begin{itemize}
        \item $(AB)C = A(BC)$
        \item $k(AB) = (kA)B = A(kB)$
        \item $A(B + C) = AB + AC; (B + C)A = BA + CA$
        \item $EA = AE = A$
        \item $(\lambda E)A = \lambda A = A(\lambda E)$
        \item $A^k A^l = A^{k+l}$
        \item $(A^k)^l = A^{kl}$
    \end{itemize}
    \textbf{注意}：矩阵的乘法一般不满足交换律，即 $AB$ 有意义，但 $BA$ 不一定有意义；即使 $AB$ 和 $BA$ 都有意义，两者也不一定相等。两个非零矩阵相乘，可能是零矩阵，从而不能从 $AB=0$ 必然推出 $A=0$ 或 $B=0$。矩阵的乘法一般不满足消去律，即不能从 $AC=BC$ 必然推出 $A=B$。
    \item \textbf{关于矩阵的转置运算的公式}
    \begin{itemize}
        \item $(A^T)^T = A$
        \item $(A + B)^T = A^T + B^T$
        \item $(kA)^T = k A^T$
        \item $(AB)^T = B^T A^T$
    \end{itemize}
    \item \textbf{关于方阵的行列式的公式}
    若 $A, B$ 是 $n$ 阶方阵，
    \begin{itemize}
        \item $|A^T| = |A|$
        \item $|\lambda A| = \lambda^n |A|$
        \item $|AB| = |A||B|$
        \item $|AB| = |BA|$
    \end{itemize}
\end{enumerate}

\subsection{几类特殊矩阵}
\begin{enumerate}
    \item \textbf{单位矩阵}：主对角线上元素都是 1，其余元素均为零的方阵称为单位矩阵，记为 $E$（或 $I$），即
    $$
    E = \begin{bmatrix}
    1 & 0 & \cdots & 0 \\
    0 & 1 & \cdots & 0 \\
    \vdots & \vdots & & \vdots \\
    0 & 0 & \cdots & 1
    \end{bmatrix}
    $$
    \item \textbf{对角矩阵}：主对角线上元素为任意常数，而主对角线外的元素均为零的矩阵。若对角矩阵的主对角线上的元素相等，则称为数量矩阵。
    \item \textbf{三角矩阵}：主对角线下方元素全为零的方阵称为上三角矩阵；主对角线上方元素全为零的方阵称为下三角矩阵；上、下三角矩阵统称为三角矩阵。
    \item \textbf{对称矩阵}：如果 $n$ 阶方阵 $A=(a_{ij})$ 满足 $a_{ij} = a_{ji} (i, j=1, 2, \cdots, n)$，即 $A^T = A$，则称 $A$ 为对称矩阵。
    \item \textbf{反对称矩阵}：如果 $n$ 阶方阵 $A=(a_{ij})$ 满足 $a_{ij} = -a_{ji} (i \neq j), a_{ii} = 0 (i, j=1, 2, \cdots, n)$，即 $A^T = -A$，则称 $A$ 为反对称矩阵。
    \item \textbf{正交矩阵}：对方阵 $A$，如果有 $A^T A = A A^T = E$，则称 $A$ 为正交矩阵。
    \item \textbf{幂零矩阵}：对方阵 $A$，如果存在正整数 $m$，使 $A^m = 0$，则称 $A$ 为幂零矩阵。
    \item \textbf{幂等矩阵}：满足 $A^2 = A$ 的方阵 $A$ 称为幂等矩阵。

    \item \textbf{对合矩阵}：满足 $A^2 = E$ 的方阵 $A$ 称为对合矩阵。
\end{enumerate}

\section{逆矩阵}

\subsection{逆矩阵的定义}
对于 $n$ 阶方阵 $A$，如果存在 $n$ 阶方阵 $B$ 使 $AB = BA = E$，则称 $A$ 是可逆的，并把 $B$ 称为 $A$ 的逆矩阵，记作 $A^{-1} = B$。

\subsection{关于逆矩阵的常用结论}
\begin{enumerate}
    \item 方阵 $A$ 可逆的充要条件是 $|A| \neq 0$。
    \item 若 $AB = E$ 或 $BA = E$，则 $B = A^{-1}$。
    \item $(A^{-1})^{-1} = A$
    \item $(kA)^{-1} = \frac{1}{k}A^{-1}$ （其中 $k \neq 0$）
    \item $(A^T)^{-1} = (A^{-1})^T$
    \item $(AB)^{-1} = B^{-1}A^{-1}$
    \item $|A^{-1}| = |A|^{-1}$
    \item 一般情况下，$(A+B)^{-1} \neq A^{-1} + B^{-1}$
    \item 可逆的上（下）三角矩阵的逆矩阵仍为上（下）三角矩阵。
\end{enumerate}

\subsection{伴随矩阵}
\textbf{定义}：设 $A=(a_{ij})$ 是 $n$ 阶方阵，行列式 $|A|$ 的各个元素 $a_{ij}$ 的代数余子式 $A_{ij}$ 所构成的如下的矩阵
$$
A^* = \begin{bmatrix}
A_{11} & A_{21} & \cdots & A_{n1} \\
A_{12} & A_{22} & \cdots & A_{n2} \\
\vdots & \vdots & & \vdots \\
A_{1n} & A_{2n} & \cdots & A_{nn}
\end{bmatrix}
$$
称为 $A$ 的伴随矩阵。

\textbf{常用结论}：
\begin{enumerate}
    \item $A A^* = A^* A = |A|E$
    \item 若 $A$ 可逆，则 $A^{-1} = \frac{1}{|A|}A^*$
    \item $A^* = |A|A^{-1}$，且 $A^*$ 也可逆（当 $|A| \neq 0$），$(A^*)^{-1} = (A^{-1})^* = \frac{1}{|A|}A$
    \item $(AB)^* = B^* A^*$
    \item $((A^*)^T) = (A^T)^*$
    \item $(kA)^* = k^{n-1} A^* \; (k \neq 0)$
    \item $|A^*| = |A|^{n-1} \; (n \geq 2)$
    \item $(A^*)^{n-2} A \; (n \geq 2)$
\end{enumerate}

\section{初等变换}

\subsection{矩阵的初等变换}
矩阵的初等行变换与初等列变换统称为初等变换。下列三种关于矩阵的变换称为矩阵的初等行（列）变换：
\begin{enumerate}
    \item 互换矩阵中两行（列）的位置（记为 $r_i \leftrightarrow r_j, c_i \leftrightarrow c_j$）；
    \item 以非零常数乘矩阵的某一行（列）（记为 $k r_i, k c_i$）；
    \item 将矩阵的某一行（列）的 $k$ 倍加到另一行（列）上去（记为 $r_i + k r_j, c_i + k c_j$）。
\end{enumerate}

\subsection{初等矩阵}
\begin{enumerate}
    \item \textbf{定义}：由单位矩阵 $E$ 经过一次初等变换得到的矩阵称为初等矩阵。
    \item 三种初等变换对应三种初等矩阵：
    \begin{enumerate}
        \item \textbf{互换两行或两列的位置 $E(i,j)$}：
        单位矩阵 $E$ 的第 $i$ 行和第 $j$ 行互换。
        性质：$|E(i,j)| = -1$， $E^{-1}(i,j) = E(i,j)$， $E^T(i,j) = E(i,j)$。
        \item \textbf{以数 $k \neq 0$ 乘某行或某列 $E(i(k))$}：
        单位矩阵 $E$ 的第 $i$ 行乘以 $k$。
        性质：$|E(i(k))| = k$， $E^{-1}(i(k)) = E(i(\frac{1}{k}))$， $E^T(i(k)) = E(i(k))$。
        \item \textbf{以数 $k$ 乘某行（列）加到另一行（列）上去 $E(i,j(k))$}：
        单位矩阵 $E$ 的第 $j$ 行的 $k$ 倍加到第 $i$ 行。
        性质：$|E(i,j(k))| = 1$， $E^{-1}(i,j(k)) = E(i,j(-k))$， $E^T(i,j(k)) = E(j,i(k))$。
    \end{enumerate}
    \item \textbf{初等变换与初等矩阵的联系}：
    设 $A$ 是 $m \times n$ 矩阵，对 $A$ 施行一次初等行变换，相当于在 $A$ 的左边乘以相应的 $m$ 阶初等矩阵；对 $A$ 施行一次初等列变换，相当于在 $A$ 的右边乘以相应的 $n$ 阶初等矩阵。
\end{enumerate}

\subsection{初等变换化简矩阵}
\begin{enumerate}
    \item \textbf{行阶梯形矩阵}：可画出一条阶梯线，线的下方全为 0，每个台阶只有一行，台阶数即是非零行的行数，阶梯线的竖线后面的第一个元素为非零元，就是非零行的第一个非零元。
    \item \textbf{行最简形矩阵}：非零行的第一个非零元为 1，且这些非零元所在列的其他元素都为 0。
    \item 对于任何矩阵 $A$，总可经过有限次初等行变换把它变为行阶梯形矩阵和行最简形矩阵。
    \item \textbf{标准型}：对于 $m \times n$ 矩阵 $A$，总可经过初等变换，把它化为 $F = \begin{bmatrix} E_r & 0 \\ 0 & 0 \end{bmatrix}_{m \times n}$，称为等价标准型。
\end{enumerate}

\subsection{矩阵等价}
\begin{enumerate}
    \item 矩阵 $A$ 经过一系列的初等变换得到矩阵 $B$，则称 $A, B$ 等价。特别地，$A$ 经过一系列初等行（列）变换得到 $B$，称 $A, B$ 行（列）等价。
    \item 方阵 $A$ 可逆的充要条件是在有限个初等矩阵 $P_1, P_2, \cdots, P_s$，使 $A = P_1 \cdots P_s$。
    \item 方阵 $A$ 可逆的充要条件是 $A$ 与单位矩阵等价 ($A \sim E$)。
    \item $m \times n$ 矩阵 $A$ 和 $B$ 等价的充要条件是存在 $m$ 阶可逆矩阵 $P$ 和 $n$ 阶可逆矩阵 $Q$，使 $PAQ = B$。
    \item $m \times n$ 矩阵 $A$ 和 $B$ 等价的充要条件是 $A$ 和 $B$ 有相同的秩。
\end{enumerate}

\subsection{初等变换求逆矩阵}
主要有三种方法：
\begin{enumerate}
    \item $[A \vdots E] \xrightarrow{\text{初等行变换}} [E \vdots A^{-1}]$
    \item $\begin{bmatrix} A \\ E \end{bmatrix} \xrightarrow{\text{初等列变换}} \begin{bmatrix} E \\ A^{-1} \end{bmatrix}$
    \item $\begin{bmatrix} A & E \\ E & 0 \end{bmatrix} \xrightarrow{\text{初等行、列变换}} \begin{bmatrix} E & C \\ B & 0 \end{bmatrix}$，则 $A^{-1} = BC$。
\end{enumerate}

\subsection{初等矩阵的推广}
\begin{enumerate}
    \item 设 $A_{m \times n} = (a_{ij})_{m \times n}$， $E_{ij}$ 为仅在 $(i,j)$ 位置为 1，其余位置为 0 的 $m$ 阶矩阵。
    $E_{ij}A$ 相当于把 $A$ 中第 $j$ 行换成第 $i$ 行元素，其他元素为 0。
    $AE_{ij}$ 相当于把 $A$ 中第 $i$ 列换成第 $j$ 列元素，其他元素为 0。
    \item 移位矩阵：
    设 $\alpha_i$ 为 $A$ 的行向量，$A = \begin{bmatrix} \alpha_1 \\ \vdots \\ \alpha_n \end{bmatrix}$。
    左乘 $\begin{bmatrix} 0 & 1 & & \\ & 0 & \ddots & \\ & & \ddots & 1 \\ & & & 0 \end{bmatrix}$ 相当于把矩阵 $A$ 的各行向上递推了一次（第一行丢失，最后一行补0）。
    类似地，右乘特定矩阵可以将列向量进行左右移位。
\end{enumerate}

\section{矩阵的秩}

\subsection{基本概念}
\begin{enumerate}
    \item \textbf{$k$ 阶子式}：在 $m \times n$ 矩阵 $A$ 中，任取 $k$ 行 $k$ 列，其交叉处的 $k^2$ 个元素按原顺序组成一个 $k$ 阶矩阵，其行列式称为 $A$ 的一个 $k$ 阶子式。
    \item \textbf{矩阵的秩}：矩阵 $A$ 的不为零的子式的最高阶数称为 $A$ 的秩，记为 $r(A)$。
\end{enumerate}

\subsection{常用公式和结论}
设 $A$ 为 $m \times n$ 矩阵，$B$ 为 $n \times l$ 矩阵（部分公式要求同型），则
\begin{enumerate}
    \item $0 \le r(A) \le \min\{m, n\}$
    \item $r(A^T) = r(A)$
    \item 若 $A \neq 0$，则 $r(A) \ge 1$
    \item $r(A \pm B) \le r(A) + r(B)$ （此时 $A, B$ 为同型矩阵）
    \item 若 $P$ 可逆，则 $r(PA) = r(A)$；若 $Q$ 可逆，则 $r(AQ) = r(A)$
    \item $r(AB) \le \min\{r(A), r(B)\}$
    \item $r(AB) \ge r(A) + r(B) - n$
    \item 若 $AB = 0$，则 $r(A) + r(B) \le n$
    \item $A$ 行满秩 $\iff r(A) = m \iff A$ 的等价标准型为 $(E_m, 0)$
    \item $A$ 列满秩 $\iff r(A) = n \iff A$ 的等价标准型为 $\begin{bmatrix} E_n \\ 0 \end{bmatrix}$
    \item 若 $A$ 是 $n$ 阶方阵，则 $r(A) = n \iff |A| \neq 0$；$r(A) < n \iff |A| = 0$
    \item 同型矩阵 $A, B$ 等价的充要条件是 $r(A) = r(B)$
    \item 设 $A^*$ 是 $n$ 阶方阵 $A$ 的伴随矩阵，则
    $$
    r(A^*) = \begin{cases}
    n, & r(A) = n \\
    1, & r(A) = n-1 \\
    0, & r(A) < n-1
    \end{cases}
    $$
\end{enumerate}

\section{分块矩阵}

\subsection{分块矩阵的定义}
将矩阵 $A$ 用若干条纵线和横线分成许多个小矩阵，每个小矩阵称为 $A$ 的子块，以子块为元素的形式上的矩阵称为分块矩阵。
例如：
$$
A = \begin{bmatrix} A_{11} & A_{12} & A_{13} \\ A_{21} & A_{22} & A_{23} \end{bmatrix}
$$
其中 $A_{ij}$ 是子块。

\subsection{分块矩阵的运算}
\begin{enumerate}
    \item \textbf{分块矩阵的加减法}：若矩阵 $A$ 与 $B$ 有相同的行数和列数，且采用相同的分块法，则
    $$
    A \pm B = \begin{bmatrix} A_{11} \pm B_{11} & \cdots & A_{1r} \pm B_{1r} \\ \vdots & & \vdots \\ A_{s1} \pm B_{s1} & \cdots & A_{sr} \pm B_{sr} \end{bmatrix}
    $$
    \item \textbf{分块矩阵的数乘}：
    $$
    \lambda A = \begin{bmatrix} \lambda A_{11} & \cdots & \lambda A_{1r} \\ \vdots & & \vdots \\ \lambda A_{s1} & \cdots & \lambda A_{sr} \end{bmatrix}
    $$
    \item \textbf{分块矩阵的乘法}：若 $A$ 为 $m \times l$ 矩阵，$B$ 为 $l \times n$ 矩阵，且 $A$ 的列分法与 $B$ 的行分法相同，则
    $$
    AB = C = \begin{bmatrix} C_{11} & \cdots & C_{1r} \\ \vdots & & \vdots \\ C_{s1} & \cdots & C_{sr} \end{bmatrix}
    $$
    其中 $C_{ij} = \sum_{k} A_{ik} B_{kj}$。
    \item \textbf{分块矩阵的转置}：
    $$
    A^T = \begin{bmatrix} A_{11}^T & \cdots & A_{s1}^T \\ \vdots & & \vdots \\ A_{1r}^T & \cdots & A_{sr}^T \end{bmatrix}
    $$
    注意子块的位置互换，且每个子块也要转置。
\end{enumerate}

\subsection{分块矩阵常用结论}
\begin{enumerate}
    \item 设 $A = \text{diag}(A_1, A_2, \cdots, A_m)$，其中 $A_i$ 都是方阵，则
    $$ |A| = |A_1| |A_2| \cdots |A_m| $$
    $$ A^k = \text{diag}(A_1^k, A_2^k, \cdots, A_m^k) $$
    \item 若 $A_i (i=1, 2, \cdots, m)$ 均为可逆矩阵，则
    $$
    A^{-1} = \begin{bmatrix} A_1^{-1} & & \\ & \ddots & \\ & & A_m^{-1} \end{bmatrix}
    $$
    \item 设 $A = \begin{bmatrix} & & A_1 \\ & \dots & \\ A_m & & \end{bmatrix}$（副对角线分块），且 $A_i$ 均为可逆矩阵，则
    $$
    A^{-1} = \begin{bmatrix} & & A_m^{-1} \\ & \dots & \\ A_1^{-1} & & \end{bmatrix}
    $$
    \item 设 $A, B$ 为 $n$ 阶方阵，则
    $$
    \begin{vmatrix} A & C \\ 0 & B \end{vmatrix} = \begin{vmatrix} A & 0 \\ 0 & B \end{vmatrix} = \begin{vmatrix} A & 0 \\ D & B \end{vmatrix} = |A| |B|
    $$
\end{enumerate}


\section{矩阵方程}

\begin{enumerate}
    \item 设 $A$ 为 $n$ 阶可逆矩阵，$B$ 为 $n \times s$ 矩阵，则矩阵方程 $AX = B$ 有唯一解 $X = A^{-1}B$.
    \item 设 $A$ 为 $n$ 阶可逆矩阵，$B$ 为 $s \times n$ 矩阵，则矩阵方程 $XA = B$ 有唯一解 $X = BA^{-1}$.
    \item 设 $A$ 为 $n$ 阶可逆矩阵，$C$ 为 $s$ 阶可逆矩阵，$B$ 为 $n \times s$ 矩阵，则矩阵方程 $AXC = B$ 有唯一解 $X = A^{-1}BC^{-1}$.
    \item \textbf{初等变换求解矩阵方程}
    \begin{itemize}
        \item 矩阵方程 $AX = B$ 可用下列方法求解：
        $$ (A \vdots B) \xrightarrow{\text{初等行变换}} (E \vdots A^{-1}B) $$
        \item 矩阵方程 $XA = B$ 可用下列方法求解：
        $$ \begin{bmatrix} A \\ B \end{bmatrix} \xrightarrow{\text{初等列变换}} \begin{bmatrix} E \\ BA^{-1} \end{bmatrix} $$
    \end{itemize}
\end{enumerate}


\newpage
\section{向量}

\subsection{向量的运算}

\begin{enumerate}
    \item \textbf{向量的定义}
    由 $n$ 个数 $a_1, a_2, \cdots, a_n$ 组成的有序数组称为 $n$ 维向量，简称向量。
    $$ \boldsymbol{\alpha} = (a_1, a_2, \cdots, a_n) $$
    称为 $n$ 维行向量，$a_i$ 移为 $\boldsymbol{\alpha}$ 的第 $i$ 个分量；
    $$ \boldsymbol{\beta} = \begin{bmatrix} b_1 \\ b_2 \\ \vdots \\ b_n \end{bmatrix} = (b_1, b_2, \cdots, b_n)^T $$
    称为 $n$ 维列向量。
    
    由定义看出，$n$ 维行（列）向量就是 $1 \times n$ ($n \times 1$) 矩阵。本书约定，所讨论的向量未指明是行向量还是列向量时，都当作列向量。
    分量全为 0 的向量称为零向量；设 $\boldsymbol{\alpha} = (a_1, a_2, \cdots, a_n)$，称 $-\boldsymbol{\alpha} = (-a_1, -a_2, \cdots, -a_n)$ 为 $\boldsymbol{\alpha}$ 的负向量。
    
    \item \textbf{向量的运算}
    \begin{enumerate}
        \item \textbf{向量的相等}：设 $\boldsymbol{\alpha} = (a_1, a_2, \cdots, a_n), \boldsymbol{\beta} = (b_1, b_2, \cdots, b_n)$，如果 $a_i = b_i (i=1, 2, \cdots, n)$，则称向量 $\boldsymbol{\alpha}$ 与 $\boldsymbol{\beta}$ 相等。
        \item \textbf{向量的加减}：设 $\boldsymbol{\alpha} = (a_1, a_2, \cdots, a_n), \boldsymbol{\beta} = (b_1, b_2, \cdots, b_n)$，则 $\boldsymbol{\alpha} \pm \boldsymbol{\beta} = (a_1 \pm b_1, a_2 \pm b_2, \cdots, a_n \pm b_n)$。
        \item \textbf{向量的数乘}：设 $\boldsymbol{\alpha} = (a_1, a_2, \cdots, a_n)$，$k$ 为常数，则 $k\boldsymbol{\alpha} = (ka_1, ka_2, \cdots, ka_n)$。
    \end{enumerate}
    
    \item \textbf{向量的运算规律}
    设 $\boldsymbol{\alpha}, \boldsymbol{\beta}, \boldsymbol{\gamma}$ 均为 $n$ 维向量，$\lambda, \mu$ 为实数，则
    \begin{enumerate}
        \item $\boldsymbol{\alpha} + \boldsymbol{\beta} = \boldsymbol{\beta} + \boldsymbol{\alpha}$
        \item $(\boldsymbol{\alpha} + \boldsymbol{\beta}) + \boldsymbol{\gamma} = \boldsymbol{\alpha} + (\boldsymbol{\beta} + \boldsymbol{\gamma})$
        \item $\boldsymbol{\alpha} + \boldsymbol{0} = \boldsymbol{\alpha}$
        \item $\boldsymbol{\alpha} + (-\boldsymbol{\alpha}) = \boldsymbol{0}$
        \item $1\boldsymbol{\alpha} = \boldsymbol{\alpha}$
        \item $\lambda(\mu\boldsymbol{\alpha}) = (\lambda\mu)\boldsymbol{\alpha}$
        \item $\lambda(\boldsymbol{\alpha} + \boldsymbol{\beta}) = \lambda\boldsymbol{\alpha} + \lambda\boldsymbol{\beta}$
        \item $(\lambda + \mu)\boldsymbol{\alpha} = \lambda\boldsymbol{\alpha} + \mu\boldsymbol{\alpha}$
    \end{enumerate}
\end{enumerate}

\subsection{向量间的线性关系}

\begin{enumerate}
    \item \textbf{基本概念}
    \begin{enumerate}
        \item \textbf{线性表示}：对于向量 $\beta, \alpha_1, \cdots, \alpha_m$，如果存在一组数 $k_1, \cdots, k_m$，使得
        $$ \beta = k_1\alpha_1 + \cdots + k_m\alpha_m $$
        成立，则称 $\beta$ 是 $\alpha_1, \cdots, \alpha_m$ 的线性组合，或称 $\beta$ 可由 $\alpha_1, \cdots, \alpha_m$ 线性表示。
        \item \textbf{线性相关与线性无关}：设 $\alpha_1, \cdots, \alpha_m$ 为一组向量，如果存在一组不全为零的数 $k_1, \cdots, k_m$，使得
        $$ k_1\alpha_1 + k_2\alpha_2 + \cdots + k_m\alpha_m = 0 $$
        成立，则称向量组 $\alpha_1, \cdots, \alpha_m$ 线性相关；当仅当 $k_1 = k_2 = \cdots = k_m = 0$ 时等式成立，则称向量组 $\alpha_1, \cdots, \alpha_m$ 线性无关。
    \end{enumerate}
    \item \textbf{常用结论}
    设 $\alpha_j = (a_{1j}, a_{2j}, \cdots, a_{mj})^T, \cdots, \alpha_m = (a_{m1}, a_{m2}, \cdots, a_{mn})^T$, $\beta = (b_1, b_2, \cdots, b_n)^T$，这里 $m \le n$。
    \begin{enumerate}
        \item $\beta$ 可由 $\alpha_1, \alpha_2, \cdots, \alpha_m$ 线性表示的充分必要条件是线性方程组 $x_1\alpha_1 + x_2\alpha_2 + \cdots + x_m\alpha_m = \beta$ 有解，即下列线性方程组有解
        $$
        \begin{cases}
        a_{11}x_1 + a_{12}x_2 + \cdots + a_{1m}x_m = b_1 \\
        a_{21}x_1 + a_{22}x_2 + \cdots + a_{2m}x_m = b_2 \\
        \cdots \cdots \cdots \cdots \cdots \cdots \\
        a_{n1}x_1 + a_{n2}x_2 + \cdots + a_{nm}x_m = b_n
        \end{cases}
        $$
        \item 令 $A = (\alpha_1, \alpha_2, \cdots, \alpha_m), B = (\alpha_1, \alpha_2, \cdots, \alpha_m, \beta)$。$\beta$ 可由 $\alpha_1, \alpha_2, \cdots, \alpha_m$ 线性表示的充分必要条件是 $r(A) = r(B)$。
        \item $\beta$ 可由 $\alpha_1, \alpha_2, \cdots, \alpha_m$ 唯一线性表示的充分必要条件是 $r(A) = r(B) = m$。
        \item $\beta$ 不能由 $\alpha_1, \alpha_2, \cdots, \alpha_m$ 线性表示的充分必要条件是 $r(A) < r(B)$。
        \item 向量组 $\alpha_1, \alpha_2, \cdots, \alpha_m$ 线性相关的充分必要条件是齐次线性方程组
        $$
        \begin{cases}
        a_{11}x_1 + a_{12}x_2 + \cdots + a_{1m}x_m = 0 \\
        a_{21}x_1 + a_{22}x_2 + \cdots + a_{2m}x_m = 0 \\
        \cdots \cdots \cdots \cdots \cdots \cdots \\
        a_{n1}x_1 + a_{n2}x_2 + \cdots + a_{nm}x_m = 0
        \end{cases}
        $$
        有非零解。且当 $m=n$ 时，其线性相关的充要条件是 $|A| = 0$。
        \item 向量组 $\alpha_1, \alpha_2, \cdots, \alpha_m$ 线性无关的充分必要条件是齐次线性方程组只有零解，且当 $m=n$ 时，其线性无关的充要条件是 $|A| \neq 0$。
        \item 向量组 $\alpha_1, \alpha_2, \cdots, \alpha_m$ 线性相关的充要条件是以 $\alpha_1, \alpha_2, \cdots, \alpha_m$ 为列向量的矩阵的秩小于向量个数 $m$。
        \item 向量组 $\alpha_1, \alpha_2, \cdots, \alpha_m$ 线性无关的充要条件是以 $\alpha_1, \alpha_2, \cdots, \alpha_m$ 为列向量的矩阵的秩等于向量个数 $m$。
        \item 向量组 $\alpha_1, \alpha_2, \cdots, \alpha_m (m \ge 2)$ 线性相关的充要条件是向量组至少有一个向量是其余向量的线性组合；向量组 $\alpha_1, \alpha_2, \cdots, \alpha_m (m \ge 2)$ 线性无关的充要条件是向量组中每一个向量都不能由其余向量线性表示。
        \item 如果向量组 $\alpha_1, \alpha_2, \cdots, \alpha_m$ 线性无关，而向量组 $\alpha_1, \alpha_2, \cdots, \alpha_m, \beta$ 线性相关，则 $\beta$ 可以由 $\alpha_1, \alpha_2, \cdots, \alpha_m$ 线性表示，且表达式唯一。
        \item 如果向量组 $\alpha_1, \alpha_2, \cdots, \alpha_m$ 可以由向量组 $\beta_1, \beta_2, \cdots, \beta_l$ 线性表示，并且 $m > l$，则向量组 $\alpha_1, \alpha_2, \cdots, \alpha_m$ 线性相关；或者说，如果向量组 $\alpha_1, \alpha_2, \cdots, \alpha_m$ 线性无关，并且可以由 $\beta_1, \beta_2, \cdots, \beta_l$ 线性表示，则 $m \le l$。
        \item 向量组 $\alpha_1, \alpha_2, \cdots, \alpha_m$ 中，如果有一个部分组线性相关，则整个向量组线性相关；如果整个向量组 $\alpha_1, \alpha_2, \cdots, \alpha_m$ 线性无关，则其任一部分组也一定线性无关。
        \item 设 $\alpha_i$ 为 $A$ 的行向量，$A = \begin{bmatrix} \alpha_1 \\ \vdots \\ \alpha_n \end{bmatrix}$。
    左乘 $\begin{bmatrix} 0 & 1 & & \\ & 0 & \ddots & \\ & & \ddots & 1 \\ & & & 0 \end{bmatrix}$ 相当于把矩阵 $A$ 的各行向上递推了一次（第一行丢失，最后一行补0）。
    类似地，右乘特定矩阵可以将列向量进行左右移位。
    \end{enumerate}
\end{enumerate}

\subsection{向量组的极大线性无关组和秩}

\begin{enumerate}
    \item \textbf{极大无关组}：设向量组 $\alpha_{i_1}, \alpha_{i_2}, \cdots, \alpha_{i_r}$ 为向量组 $\alpha_1, \alpha_2, \cdots, \alpha_m$ 的一个部分组，且满足
    \begin{enumerate}
        \item $\alpha_{i_1}, \alpha_{i_2}, \cdots, \alpha_{i_r}$ 线性无关，
        \item 向量组 $\alpha_1, \alpha_2, \cdots, \alpha_m$ 中任一向量均可由 $\alpha_{i_1}, \alpha_{i_2}, \cdots, \alpha_{i_r}$ 线性表示，
    \end{enumerate}
    则称向量组 $\alpha_{i_1}, \alpha_{i_2}, \cdots, \alpha_{i_r}$ 为向量组 $\alpha_1, \alpha_2, \cdots, \alpha_m$ 的一个极大线性无关组，简称极大无关组。

    \item \textbf{向量组的秩}：向量组 $\alpha_1, \alpha_2, \cdots, \alpha_m$ 的极大无关组中所含向量的个数称为该向量组的秩，记为 $r(\alpha_1, \alpha_2, \cdots, \alpha_m)$。
    如果一个向量组仅含有零向量，则规定它的秩为零。

    \item \textbf{向量组的秩的性质}
    \begin{enumerate}
        \item 若 $r(\alpha_1, \alpha_2, \cdots, \alpha_m) = r$，则
        \begin{itemize}
            \item $\alpha_1, \cdots, \alpha_m$ 的任何含有多于 $r$ 个向量的部分组一定线性相关；
            \item $\alpha_1, \cdots, \alpha_m$ 的任何含 $r$ 个向量的线性无关部分组一定是极大无关组。
        \end{itemize}
        \item $r(\alpha_1, \alpha_2, \cdots, \alpha_m) \le m$，且 $r(\alpha_1, \alpha_2, \cdots, \alpha_m) = m \iff \alpha_1, \cdots, \alpha_m$ 线性无关。
        \item 向量 $\beta$ 可用 $\alpha_1, \alpha_2, \cdots, \alpha_m$ 线性表示 $\iff r(\alpha_1, \cdots, \alpha_m) = r(\alpha_1, \cdots, \alpha_m, \beta)$。
        \item 若 $\beta_1, \beta_2, \cdots, \beta_s$ 可用 $\alpha_1, \alpha_2, \cdots, \alpha_m$ 线性表示，则 $r(\beta_1, \cdots, r(\beta_s) \le r(\alpha_1, \cdots, \alpha_m)$。
        \item 设 $A$ 是一个 $m \times n$ 矩阵，记 $\alpha_1, \alpha_2, \cdots, \alpha_n$ 是 $A$ 的列向量组 ($m$ 维)，$\beta_1, \beta_2, \cdots, \beta_m$ 是 $A$ 的行向量组 ($n$ 维)，则 $r(A) = r(\alpha_1, \alpha_2, \cdots, \alpha_n) = r(\beta_1, \beta_2, \cdots, \beta_m)$。
    \end{enumerate}

    \item \textbf{向量组的等价}：两个向量组能够相互线性表示，则称这两个向量组等价。
    \begin{itemize}
        \item 任一向量组和它的极大无关组等价。
        \item 向量组的任意两个极大无关组等价。
        \item 具有相同秩的向量组一定等价。
    \end{itemize}
\end{enumerate}

\subsection{向量的内积与向量空间}

\begin{enumerate}
    \item \textbf{向量的内积}：给定 $R^n$ 中向量 $\alpha=(a_1, a_2, \cdots, a_n)^T, \beta=(b_1, b_2, \cdots, b_n)^T$，则称 $\sum_{i=1}^n a_i b_i$ 为向量 $\alpha$ 与 $\beta$ 的内积，记为 $(\alpha, \beta)$，即 $(\alpha, \beta) = \alpha^T \beta = \sum_{i=1}^n a_i b_i$。
    内积具有下列性质：
    \begin{enumerate}
        \item $(\alpha, \beta) = (\beta, \alpha)$;
        \item $(k\alpha, \beta) = k(\alpha, \beta)$;
        \item $(\alpha + \beta, \gamma) = (\alpha, \gamma) + (\beta, \gamma)$;
        \item $(\alpha, \alpha) \ge 0$，当且仅当 $\alpha=0$ 时，等号成立。
    \end{enumerate}

    \item \textbf{向量的范数}：设 $\alpha$ 为 $R^n$ 中任意向量，将非负实数 $\sqrt{(\alpha, \alpha)}$ 定义为 $\alpha$ 的长度，记为 $||\alpha||$，即若 $\alpha=(a_1, a_2, \cdots, a_n)^T$，则有 $||\alpha|| = \sqrt{a_1^2 + a_2^2 + \cdots + a_n^2}$。
    向量的长度也称为向量的范数或模。
    向量范数具有下列性质：
    \begin{enumerate}
        \item $||\alpha|| \ge 0$，当且仅当 $\alpha=0$ 时，等号成立；
        \item 对于任意向量 $\alpha$ 和任意实数 $k$，都有 $||k\alpha|| = |k| ||\alpha||$;
        \item 对于任意 $n$ 维向量 $\alpha$ 和 $\beta$，$|(\alpha, \beta)| \le ||\alpha|| ||\beta||$。
    \end{enumerate}

    \item \textbf{向量的正交}：如果向量 $\alpha$ 和 $\beta$ 的内积等于零，即 $(\alpha, \beta)=0$，则称 $\alpha$ 和 $\beta$ 相互正交。
    如果非零向量组 $\alpha_1, \cdots, \alpha_s$ 中向量两两正交，即 $(\alpha_i, \alpha_j)=0 (i \neq j, i,j=1,2,\cdots,s)$，则称该向量组为正交向量组。
    正交向量具有下列性质：
    \begin{enumerate}
        \item 零向量与任何向量正交；
        \item 与自己正交的向量只有零向量；
        \item 正交向量组是线性无关的。
    \end{enumerate}

    \item 对于任意向量 $\alpha$ 和 $\beta$，有三角不等式 $||\alpha + \beta|| \le ||\alpha|| + ||\beta||$。
    当且仅当 $\alpha$ 与 $\beta$ 相互正交时，有 $||\alpha + \beta||^2 = ||\alpha||^2 + ||\beta||^2$。

    \item \textbf{向量空间}：设 $V$ 是实数域 $R$ 上的 $n$ 维向量组成的集合，如果 $V$ 关于向量的加法和数乘是封闭的，即若 $\alpha \in V, \beta \in V$，则 $\alpha + \beta \in V$；若 $\alpha \in V, k \in R$，则 $k\alpha \in V$，则称 $V$ 是实数域 $R$ 上的向量空间。
    显然，实数域 $R$ 上的 $n$ 维向量的全体构成一个向量空间，记为 $R^n$。

    \item \textbf{基与坐标}：在向量空间 $V$ 中一个线性无关的向量组 $\xi_1, \cdots, \xi_n$ 称为 $V$ 的一组基。若 $\alpha \in R^n$ 为任一向量，且 $\alpha = a_1\xi_1 + \cdots + a_n\xi_n$，则称 $a_1, \cdots, a_n$ 为 $\alpha$ 关于基 $\xi_1, \cdots, \xi_n$ 的坐标，记作 $\alpha = (a_1, \cdots, a_n)^T$。

    \item \textbf{基变换与坐标变换}：设 $\xi_1, \cdots, \xi_n$ 和 $\eta_1, \cdots, \eta_n$ 是 $R^n$ 的两组基，且有
    $$ (\eta_1, \cdots, \eta_n) = (\xi_1, \cdots, \xi_n) A $$
    其中 $A = \begin{pmatrix} a_{11} & \cdots & a_{1n} \\ \vdots & & \vdots \\ a_{n1} & \cdots & a_{nn} \end{pmatrix}$。
    称 $A$ 为由基 $\xi_1, \cdots, \xi_n$ 到基 $\eta_1, \cdots, \eta_n$ 的过渡矩阵。两个基之间的过渡矩阵是可逆矩阵。
    设 $\alpha \in R^n$ 在基 $\xi_1, \cdots, \xi_n$ 和基 $\eta_1, \cdots, \eta_n$ 下的坐标分别为 $(x_1, x_2, \cdots, x_n)^T$ 与 $(y_1, y_2, \cdots, y_n)^T$，則有
    $$ \begin{pmatrix} x_1 \\ \vdots \\ x_n \end{pmatrix} = A \begin{pmatrix} y_1 \\ \vdots \\ y_n \end{pmatrix} \quad \text{或} \quad \begin{pmatrix} y_1 \\ \vdots \\ y_n \end{pmatrix} = A^{-1} \begin{pmatrix} x_1 \\ \vdots \\ x_n \end{pmatrix} $$
    称为坐标变换公式。

    \item \textbf{$R^n$ 的标准正交基}：向量空间 $R^n$ 中一个向量组 $\eta_1, \cdots, \eta_n$ 满足
    \begin{enumerate}
        \item 两两正交，即 $(\eta_i, \eta_j)=0, i \neq j$；
        \item 都是单位向量，即 $||\eta_i||=1, i=1, \cdots, n$。
    \end{enumerate}
    则称 $\eta_1, \cdots, \eta_n$ 为 $R^n$ 的一组标准正交基。

    \item \textbf{标准正交基的求法}
    \begin{enumerate}
        \item \textbf{施密特正交化方法}：给定一线性无关向量组 $\alpha_1, \cdots, \alpha_s$，由其生成等价的 $s$ 个向量的正交向量组 $\beta_1, \cdots, \beta_s$ 的公式如下：
        $$ \beta_1 = \alpha_1 $$
        $$ \beta_2 = \alpha_2 - \frac{(\alpha_2, \beta_1)}{(\beta_1, \beta_1)}\beta_1 $$
        $$ \beta_s = \alpha_s - \frac{(\alpha_s, \beta_1)}{(\beta_1, \beta_1)}\beta_1 - \frac{(\alpha_s, \beta_2)}{(\beta_2, \beta_2)}\beta_2 - \cdots - \frac{(\alpha_s, \beta_{s-1})}{(\beta_{s-1}, \beta_{s-1})}\beta_{s-1} $$
        \item \textbf{给定 $R^n$ 任意一组基，把它变为标准正交基的步骤如下}：
        \begin{enumerate}
            \item 利用施密特正交化方法，由这组基生成 $n$ 个向量的正交向量组；
            \item 把正交向量组中每个向量标准化，即单位化。
        \end{enumerate}
        这样就得到 $R^n$ 的一组标准正交基，这一过程称为标准正交化。
    \end{enumerate}

    \item \textbf{两组标准正交基之间的过渡矩阵}：设 $R^n$ 的两组标准正交基 $\xi_1, \xi_2, \cdots, \xi_n$ 和 $\eta_1, \eta_2, \cdots, \eta_n$ 间的过渡矩阵为 $Q$，则存在下列关系
    $$ (\xi_1, \xi_2, \cdots, \xi_n) = (\eta_1, \eta_2, \cdots, \eta_n) Q $$
    且 $Q$ 满足 $Q^T Q = E$，即 $Q$ 为正交矩阵。
\end{enumerate}

\newpage
\section{矩阵的特征值与特征向量}

\subsection{基本概念}
\begin{enumerate}
    \item \textbf{特征值与特征向量}：设 $A$ 为 $n$ 阶方阵，若存在一个数 $\lambda$ 及非零的 $n$ 维列向量 $x$ 使得 $Ax = \lambda x$ 成立，则称 $\lambda$ 为 $A$ 的一个特征值，称非零向量 $x$ 为属于特征值 $\lambda$ 的特征向量。
    \item \textbf{特征矩阵、特征多项式、特征方程}：$A-\lambda E$ 称为 $A$ 的特征矩阵， $|A-\lambda E|$ 称为 $A$ 的特征多项式， $|A-\lambda E|=0$ 称为 $A$ 的特征方程。
\end{enumerate}

\subsection{特征值的性质及运算}
若 $\lambda$ 是 $A$ 的特征值，则
\begin{enumerate}
    \item $k\lambda$ 是 $kA$ 的特征值。
    \item $\lambda^m$ 是 $A^m$ 的特征值。
    \item $f(\lambda) = \sum_{i=0}^m a_i \lambda^i$ 是 $f(A) = \sum_{i=0}^m a_i A^i$ 的特征值。
    \item 若 $A$ 可逆，则 $\lambda \neq 0$，且 $\frac{1}{\lambda}$ 是 $A^{-1}$ 的特征值。
    \item 若 $\lambda \neq 0$，则 $\frac{|A|}{\lambda}$ 是 $A^*$ 的特征值。
    \item $A$ 与 $A^T$ 有相同的特征值。
    \item $AB$ 与 $BA$ 有相同的特征值。
    \item $A$ 是奇异矩阵（$|A|=0$）的充分必要条件是 $A$ 有特征值 0。
    \item 幂零矩阵只有 0 为特征值。
    \item 单位矩阵 $E$ 有 $n$ 重特征值 1。
    \item 数量矩阵 $kE$ 有 $n$ 重特征值 $k$。
    \item 幂零矩阵 ($A^k=0$) 有 $n$ 重特征值 0。
    \item 幂等矩阵 ($A^2=A$) 的特征值只能是 0 或 1。
    \item 对合矩阵 ($A^2=E$) 的特征值只能是 1 或 -1。
    \item 若 $A$ 的 $n$ 个特征值为 $\lambda_1, \lambda_2, \cdots, \lambda_n$，则
    \begin{enumerate}
        \item $\lambda_1 + \lambda_2 + \cdots + \lambda_n = a_{11} + a_{22} + \cdots + a_{nn}$，即特征值之和等于矩阵的迹。
        \item $\lambda_1 \lambda_2 \cdots \lambda_n = |A|$，即特征值之积等于矩阵的行列式。
    \end{enumerate}
\end{enumerate}

\subsection{特征向量的性质}
\begin{enumerate}
    \item 若 $x$ 是 $A$ 的属于特征值 $\lambda$ 的特征向量，则 $kx$ ($k \neq 0$) 一定是非零向量，且也是 $A$ 的属于特征值 $\lambda$ 的特征向量。
    \item 若 $x_1, x_2, \cdots, x_m$ 都是 $A$ 的属于同一特征值 $\lambda$ 的特征向量，且 $k_1 x_1 + k_2 x_2 + \cdots + k_m x_m \neq 0$，则 $k_1 x_1 + k_2 x_2 + \cdots + k_m x_m$ 也是 $A$ 的属于特征值 $\lambda$ 的特征向量。
    \item 设 $\lambda_1, \lambda_2$ 是 $A$ 的两个不同特征值，对应的特征向量为 $x_1, x_2$，则 $x_1+x_2$ 不是 $A$ 的特征向量。
    \item 若 $\lambda$ 是 $A$ 的 $r$ 重特征值，则属于 $\lambda$ 的线性无关的特征向量最多有 $r$ 个。
    \item $A$ 的属于不同特征值的特征向量线性无关。
    \item 设 $\lambda$ 是 $A$ 的特征值，$x$ 是属于 $\lambda$ 的特征向量，则
    \begin{enumerate}
        \item $x$ 是 $kA$ 的属于特征值 $k\lambda$ 的特征向量，$A^m$ 的属于特征值 $\lambda^m$ 的特征向量，和 $f(A)=\sum a_i A^i$ 的属于特征值 $f(\lambda)=\sum a_i \lambda^i$ 的特征向量。
        \item 若 $\lambda \neq 0$，则 $x$ 也是 $A^{-1}$ 的属于特征值 $\frac{1}{\lambda}$ 的特征向量和 $A^*$ 的属于特征值 $\frac{|A|}{\lambda}$ 的特征向量。
        \item 若 $B=P^{-1}AP$，$P$ 可逆，则 $\lambda$ 是 $B$ 的特征值，且 $P^{-1}x$ 是 $B$ 的属于特征值 $\lambda$ 的特征向量。
    \end{enumerate}
\end{enumerate}

\subsection{矩阵的特征值和特征向量的求法}
\begin{enumerate}
    \item 对于数字型矩阵 $A$，其特征值和特征向量的求法如下：
    \begin{enumerate}
        \item 计算特征多项式 $|A - \lambda E|$。
        \item 求解特征方程 $|A - \lambda E| = 0$，其所有根 $\lambda_1, \lambda_2, \cdots, \lambda_n$ 即为 $A$ 的全部特征值。
        \item 固定一个特征值 $\lambda_i$，求解齐次线性方程组 $(A - \lambda_i E)x = 0$，求得基础解系 $\eta_1, \eta_2, \cdots, \eta_s$，则 $A$ 的属于特征值 $\lambda_i$ 的所有特征向量为 $k_1 \eta_1 + k_2 \eta_2 + \cdots + k_s \eta_s$，其中 $k_1, k_2, \cdots, k_s$ 不全为零。
    \end{enumerate}
    \item 对于抽象型矩阵 $A$，其特征值和特征向量的求法通常有两种思路：
    \begin{enumerate}
        \item 利用特征值和特征向量的定义，若数 $\lambda$ 及非零向量 $x$ 满足 $Ax=\lambda x$，则 $\lambda$ 为 $A$ 的特征值，$x$ 为属于 $\lambda$ 的特征向量。
        \item 利用特征值和特征向量的性质，例如已知 $A$ 的特征值，便可求得 $kA, A^m, \sum a_i A^i$ 等矩阵的特征值。
    \end{enumerate}
\end{enumerate}

\newpage
\section{实对称矩阵与二次型}

\subsection{实对称矩阵的正交相似对角化}
\begin{enumerate}
    \item \textbf{实对称矩阵的性质}
    设 $A$ 是实对称矩阵，则
    \begin{enumerate}
        \item $A$ 的特征值为实数，特征向量为实向量。
        \item $A$ 的不同特征值所对应的特征向量正交。
        \item $A$ 的 $k$ 重特征值恰有 $k$ 个线性无关的特征向量。
        \item $A$ 相似于对角矩阵，且存在正交矩阵 $P$，使得 $P^{-1}AP = P^T AP = \Lambda = \text{diag}(\lambda_1, \lambda_2, \cdots, \lambda_n)$，其中 $\lambda_1, \cdots, \lambda_n$ 为 $A$ 的特征值。
        \item 实对称矩阵 $A$ 与 $B$ 相似的充要条件是它们有相同的特征值。
    \end{enumerate}
    \item \textbf{实对称矩阵正交相似对角化的步骤}
    \begin{enumerate}
        \item 求出矩阵 $A$ 的全部特征值 $\lambda_1, \lambda_2, \cdots, \lambda_n$。
        \item 对每个特征值 $\lambda_i$，求出齐次线性方程组 $(\lambda_i E - A)x=0$ 的基础解系，得到线性无关的特征向量。
        \item 将属于同一特征值的特征向量进行施密特正交化和单位化，得到两两正交的单位特征向量。
        \item 将这 $n$ 个两两正交的单位特征向量作为列向量构成正交矩阵 $P$，则有 $P^{-1}AP = P^T AP = \Lambda$。注意 $\Lambda$ 中对角元素的排列次序应与 $P$ 中列向量的排列次序相对应。
    \end{enumerate}
\end{enumerate}

\subsection{二次型}
\begin{enumerate}
    \item \textbf{二次型的基本概念}
    \begin{enumerate}
        \item \textbf{定义}：含有 $n$ 个变量 $x_1, \cdots, x_n$ 的二次齐次函数 $f(x_1, \cdots, x_n) = \sum_{i=1}^n \sum_{j=1}^n a_{ij} x_i x_j$（其中 $a_{ij}=a_{ji}$）称为 $n$ 元二次型。
        \item \textbf{矩阵表示}：二次型可表示为 $f(x) = x^T A x$，其中 $x=(x_1, \cdots, x_n)^T$，$A=(a_{ij})$ 是对称矩阵，称为二次型的矩阵。二次型的秩定义为矩阵 $A$ 的秩 $r(A)$。
        \item \textbf{标准形}：只含平方项的二次型称为标准形，即 $f = d_1 y_1^2 + d_2 y_2^2 + \cdots + d_n y_n^2$。
        \item \textbf{规范形}：在标准形中，系数仅为 $1, -1, 0$ 的二次型称为规范形，即 $f = y_1^2 + \cdots + y_p^2 - y_{p+1}^2 - \cdots - y_r^2$。
        \item \textbf{正定性}：
        \begin{itemize}
            \item \textbf{正惯性指数} $p$：标准形中正系数的个数。
            \item \textbf{负惯性指数} $q$：标准形中负系数的个数。
            \item \textbf{符号差} $p-q$。
        \end{itemize}
    \end{enumerate}
    \item \textbf{二次型的常用结论}
    \begin{enumerate}
        \item 任意二次型 $x^T A x$ 总可以通过可逆线性变换 $x=Cy$ 化为标准形 $d_1 y_1^2 + \cdots + d_n y_n^2$。若使用正交变换 $x=Py$（$P$为正交矩阵），则标准形的系数 $d_i$ 就是矩阵 $A$ 的特征值 $\lambda_i$。
        \item \textbf{惯性定理}：任意二次型 $x^T A x$ 可通过可逆线性变换化为唯一的规范形。其中的正惯性指数 $p$ 和负惯性指数 $q$ 是唯一的。
        \item 可逆线性变换不改变二次型的正定性。
    \end{enumerate}
\end{enumerate}

\subsection{二次型的正定性}
\begin{enumerate}
    \item \textbf{基本概念}：如果二次型 $f(x_1, \cdots, x_n)$ 对任意一组不全为零的实数 $x_1, \cdots, x_n$ 都有 $f(x_1, \cdots, x_n) > 0$，则称该二次型为正定二次型，正定二次型的矩阵称为正定矩阵。
    \item \textbf{实对称矩阵正定性的判定}：设 $A$ 为 $n$ 阶实对称矩阵，则下列命题等价：
    \begin{enumerate}
        \item $A$ 是正定矩阵。
        \item $A$ 的正惯性指数 $p=n$。
        \item $A$ 的顺序主子式全大于 0。
        \item $A$ 的所有主子式全大于 0。
        \item $A$ 合同于单位矩阵 $E$。
        \item $A$ 的特征值全大于 0。
        \item 存在可逆矩阵 $P$，使 $A=P^T P$。
        \item 存在非退化的上（下）三角矩阵 $Q$，使 $A=Q^T Q$。
    \end{enumerate}
    \item \textbf{正定矩阵的性质}
    \begin{enumerate}
        \item 若 $A$ 为正定矩阵，则 $|A|>0$，$A$ 为可逆对称矩阵。
        \item 若 $A$ 为正定矩阵，则 $A$ 的主对角元素 $a_{ii} > 0$。
        \item 若 $A$ 为正定矩阵，则 $A^{-1}, kA (k>0)$ 均为正定矩阵。
        \item 若 $A$ 为正定矩阵，则 $A^*, A^m$ ($m$ 为正整数) 均为正定矩阵。
        \item 若 $A, B$ 为 $n$ 阶正定矩阵，则 $A+B$ 是正定矩阵。
    \end{enumerate}
\end{enumerate}

\subsection{矩阵的合同}
\begin{enumerate}
    \item \textbf{矩阵合同的定义}：设 $A, B$ 为两个方阵，若存在可逆矩阵 $Q$，使 $B=Q^T A Q$ 成立，则称 $A$ 与 $B$ 合同。
    \item \textbf{矩阵合同的性质}
    \begin{enumerate}
        \item 矩阵 $A$ 与 $B$ 合同的充要条件是对 $A$ 的行和列施以相同的初等变换变成 $B$。
        \item 矩阵 $A$ 与 $B$ 合同的必要条件是 $A$ 与 $B$ 的秩相同。
        \item 设 $A$ 与 $B$ 是实对称矩阵，则 $A$ 与 $B$ 合同的充要条件是二次型 $x^T A x$ 与 $x^T B x$ 有相同的正负惯性指数。
        \item $A$ 与 $B$ 合同的充分条件是 $A$ 与 $B$ 相似。
    \end{enumerate}
\end{enumerate}
\newpage
\section{八大常见类型的行列式及其解法}

\begin{mybox}{八大类型行列式概览}
常见的行列式类型及其解法技巧如下：
\end{mybox}

\subsection{1. 三角形行列式}
\begin{mybox}{结构与解法}
主对角线以上（或以下）元素全为零，值为主对角线元素之积：
\(
\begin{vmatrix}
 a_{11} & * & * \\
 0 & a_{22} & * \\
 0 & 0 & a_{33}
\end{vmatrix} = a_{11} a_{22} a_{33}
\)
\end{mybox}

\subsection{2. 对角行列式}
\begin{mybox}{结构与解法}
仅主对角线有元素，其余为零，值为主对角线元素之积。
\end{mybox}

\subsection{3. 斜对角（副对角线）行列式}
\begin{mybox}{结构与解法}
仅副对角线有元素，其余为零，值为副对角线元素之积，符号为 $(-1)^{n(n-1)/2}$。
\end{mybox}

\subsection{4. 范德蒙行列式}
\begin{mybox}{结构与解法}
\(
\begin{vmatrix}
1 & 1 & \cdots & 1 \\
 a_1 & a_2 & \cdots & a_n \\
 a_1^2 & a_2^2 & \cdots & a_n^2 \\
 \vdots & \vdots & & \vdots \\
 a_1^{n-1} & a_2^{n-1} & \cdots & a_n^{n-1}
\end{vmatrix} = \prod_{1 \leq j < i \leq n} (a_i - a_j)
\)
\end{mybox}

\subsection{5. 递推型行列式}
\begin{mybox}{结构与解法}
每一行（或列）元素与前一行（或列）递推相关，常用初等变换或归纳法化简。
\end{mybox}

\subsection{6. 分块型行列式}
\begin{mybox}{结构与解法}
将行列式分为若干子块，利用分块矩阵性质或按块展开。
\end{mybox}

\subsection{7. 对称型行列式}
\begin{mybox}{结构与解法}
元素关于主对角线对称，常用对称性化简或配方法。
\end{mybox}

\subsection{8. 循环型行列式}
\begin{mybox}{结构与解法}
各行（或列）元素循环排列，常用特征根法或傅里叶变换法。
\end{mybox}

\begin{mybox}{解题建议}
遇到复杂行列式时，优先观察结构类型，选择对应技巧（如初等变换、按行列展开、分块、递归、配方法等），可大幅简化计算。
\end{mybox}
\end{document}